\documentclass[UTF8,11pt]{ctexart}
\usepackage[a4paper,margin=24mm]{geometry}
\usepackage{amsmath,amssymb,amsthm,mathtools,mathrsfs}
\usepackage[all]{xy}
\usepackage{xurl,hyperref,enumitem}
\usepackage{tikz}
\usetikzlibrary{arrows.meta,cd,decorations.markings,calc,patterns}
\hypersetup{hidelinks,breaklinks=true}
\setlength{\emergencystretch}{3em}
\allowdisplaybreaks
\newtheorem*{theorem}{Theorem}
\newtheorem*{thm}{Theorem}
\newtheorem*{lemma}{Lemma}
\newtheorem*{proposition}{Proposition}
\newtheorem*{prop}{Proposition}
\newtheorem*{corollary}{Corollary}
\newtheorem*{cor}{Corollary}
\newtheorem*{claim}{Claim}
\theoremstyle{definition}
\newtheorem*{definition}{Definition}
\newtheorem*{defn}{Definition}
\newtheorem*{remark}{Remark}
\newtheorem*{example}{Example}
\newcommand{\R}{\mathbb R}
\newcommand{\C}{\mathbb C}
\newcommand{\Z}{\mathbb Z}
\newcommand{\Q}{\mathbb Q}
\newcommand{\N}{\mathbb N}
\newcommand{\F}{\mathbb F}
\newcommand{\D}{\mathbb D}
\newcommand{\bB}{\mathbb B}
\newcommand{\bC}{\mathbb C}
\newcommand{\bR}{\mathbb R}
\newcommand{\bZ}{\mathbb Z}
\newcommand{\bN}{\mathbb N}
\newcommand{\bQ}{\mathbb Q}
\newcommand{\bD}{\mathbb D}
\newcommand{\bF}{\mathbb F}
\newcommand{\CP}{\mathbb{CP}}
\newcommand{\RP}{\mathbb{RP}}
\newcommand{\sO}{\mathscr O}
\newcommand{\sI}{\mathscr I}
\newcommand{\sF}{\mathscr F}
\newcommand{\sG}{\mathscr G}
\newcommand{\sH}{\mathscr H}
\newcommand{\sR}{\mathscr R}
\newcommand{\sM}{\mathscr M}
\newcommand{\sV}{\mathscr V}
\newcommand{\sU}{\mathscr U}
\DeclareMathOperator{\Hom}{Hom}
\DeclareMathOperator{\Ext}{Ext}
\DeclareMathOperator{\Tor}{Tor}
\DeclareMathOperator{\Spec}{Spec}
\DeclareMathOperator{\Proj}{Proj}
\DeclareMathOperator{\supp}{supp}
\DeclareMathOperator{\im}{im}
\DeclareMathOperator{\id}{id}
\DeclareMathOperator{\Tr}{Tr}
\DeclareMathOperator{\tr}{tr}
\DeclareMathOperator{\rank}{rank}
\DeclareMathOperator{\ord}{ord}
\DeclareMathOperator{\dist}{dist}
\DeclareMathOperator{\End}{End}
\DeclareMathOperator{\Aut}{Aut}
\DeclareMathOperator{\Gal}{Gal}
\DeclareMathOperator{\vol}{vol}
\DeclareMathOperator{\Cl}{Cl}
\DeclareMathOperator{\coker}{coker}
\DeclareMathOperator{\codim}{codim}
\DeclareMathOperator{\divergence}{div}
\DeclareMathOperator{\Ric}{Ric}
\DeclareMathOperator{\grad}{grad}
\DeclareMathOperator{\Hess}{Hess}
\newcommand{\abs}[1]{\left\lvert#1\right\rvert}
\newcommand{\sabs}[1]{\lvert#1\rvert}
\newcommand{\babs}[1]{\bigl\lvert#1\bigr\rvert}
\newcommand{\Babs}[1]{\Bigl\lvert#1\Bigr\rvert}
\newcommand{\bbabs}[1]{\biggl\lvert#1\biggr\rvert}
\newcommand{\BBabs}[1]{\Biggl\lvert#1\Biggr\rvert}
\newcommand{\norm}[1]{\left\lVert#1\right\rVert}
\newcommand{\snorm}[1]{\lVert#1\rVert}
\newcommand{\bnorm}[1]{\bigl\lVert#1\bigr\rVert}
\newcommand{\Bnorm}[1]{\Bigl\lVert#1\Bigr\rVert}
\newcommand{\bbnorm}[1]{\biggl\lVert#1\biggr\rVert}
\newcommand{\BBnorm}[1]{\Biggl\lVert#1\Biggr\rVert}
\newcommand{\widebar}[1]{\overline{#1}}
\newcommand{\nicefrac}[2]{\frac{#1}{#2}}
\newcommand{\linnprod}[2]{\langle#1,#2\rangle}
\newcommand{\myindex}[1]{#1}
\newcommand{\glsadd}[1]{}
\newcommand{\avoidbreak}{}
\newcommand{\sourceref}[1]{\textnormal{[原文：\nolinkurl{#1}]}}
\newcommand{\thmref}[1]{\sourceref{#1}}
\newcommand{\lemmaref}[1]{\sourceref{#1}}
\newcommand{\propref}[1]{\sourceref{#1}}
\newcommand{\corref}[1]{\sourceref{#1}}
\newcommand{\defnref}[1]{\sourceref{#1}}
\newcommand{\exampleref}[1]{\sourceref{#1}}
\newcommand{\exerciseref}[1]{\sourceref{#1}}
\newcommand{\figureref}[1]{\sourceref{#1}}
\newcommand{\sectionref}[1]{\sourceref{#1}}
\newcommand{\appendixref}[1]{\sourceref{#1}}
\newcommand{\stacksref}[2]{\href{https://stacks.math.columbia.edu/tag/#1}{#2 [#1]}}


\newcommand{\E}{\mathbb E}
\newcommand{\Prob}{\mathbb P}
\newcommand{\Reff}{\mathcal R_{\mathrm{eff}}}
\newcommand{\eps}{\varepsilon}
\DeclareMathOperator{\diam}{diam}
\DeclareMathOperator{\Var}{Var}
\DeclareMathOperator{\Crit}{Crit}
\newcommand{\dd}{\,\mathrm d}

\begin{document}
\section*{055\quad 小倍增集合的迭代和差集界}
\subsection*{来源与计数目标}
MIT 18.217, \emph{Graph Theory and Additive Combinatorics}, Fall 2019, Chapter 7, \S7.2, Theorems 7.18, 7.22, 7.24 and Lemma 7.25, printed pp.~144--146. Lectures by Yufei Zhao; lecture notes written by the students and edited with his help, as credited by the course. 取得日期：2026-09-28。
\par\url{https://ocw.mit.edu/courses/18-217-graph-theory-and-additive-combinatorics-fall-2019/c145c88ab6b56bf04e52b1bc51a2d592_MIT18_217F19_ch7.pdf}
\par\textbf{本文件只计一个问题：}从一次和集界取得任意迭代和差集的 Pl"unnecke--Ruzsa 界；7.18 和 7.25 是同一目标的支撑，不另计题数。
\subsection*{作者命题与人类证明}
\paragraph{Notation (from the source).}
\noindent\textit{以下背景与记号据所列原文章节摘编；不是逐字引文，也不是新增证明。}\par For finite subsets of an abelian group, $A+B=\{a+b:a\in A,b\in B\}$ and $A-B=\{a-b:a\in A,b\in B\}$. The notation $mA$ denotes the sum of $m$ copies of $A$.
\begin{theorem}[7.18, Ruzsa triangle inequality]
If $A,B,C$ are finite subsets of an abelian group, then
\[|A||B-C|\leq |A-B||A-C|.\]
\end{theorem}
\begin{proof}
We will construct an injection
\[\varphi:A\times(B-C)\hookrightarrow(A-B)\times(A-C).\]
For each $d\in B-C$, we can choose $b(d)\in B,c(d)\in C$ such that $d=b(d)-c(d)$. Then define $\varphi(a,d)=(a-b(d),a-c(d))$. This is injective because if $\varphi(a,d)=(x,y)$, then we can recover $(a,d)$ from $(x,y)$ because $d=y-x$ and $a=x+b(y-x)$.
\end{proof}
\begin{theorem}[7.22, Pl\"unnecke--Ruzsa inequality]
If $A$ is a finite subset of an abelian group and $|A+A|\leq K|A|$, then $|mA-nA|\leq K^{m+n}|A|$.
\end{theorem}
In proving this theorem, we will generalize to the following theorem. Set $B=A$ to recover Theorem 7.22.
\begin{theorem}[7.24]
If $A$ and $B$ are finite subsets of an abelian group and $|A+B|\leq K|A|$, then $|mB-nB|\leq K^{m+n}|A|$.
\end{theorem}
Petridis' proof relies on the following key lemma.
\begin{lemma}[7.25]
Suppose $A$ and $B$ are finite subsets of an abelian group. If $X\subseteq A$ is a nonempty subset which minimizes $|X+B|/|X|$, and $K'=|X+B|/|X|$, then $|X+B+C|\leq K'|X+C|$ for all finite sets $C$.
\end{lemma}
\begin{proof}[Proof of Theorem 7.24 assuming Lemma 7.25]
Assuming the key lemma, let us prove the theorem. Let $X$ be a nonempty subset of $A$ minimizing $|X+B|/|X|$, and let $K'=|X+B|/|X|$. Note that $K'\leq K$ by minimality. Applying the lemma with $C=rB$ where $r\geq1$, we have
\[|X+(r+1)B|\leq K'|X+rB|\leq K|X+rB|,\]
so by induction $|X+rB|\leq K^r|X|$ for all $r\geq0$. Applying Theorem 7.18 we have
\[|mB-nB|\leq\frac{|X+mB||X+nB|}{|X|}\leq K^{m+n}|X|\leq K^{m+n}|A|.\]
\end{proof}
\begin{proof}[Proof of Lemma 7.25]
We will proceed by induction on $|C|$. The base case of $|C|=1$ is clear because for any finite set $S$, $S+C$ is a translation of $S$ so $|S+C|=|S|$, thus $|X+B+C|=|X+B|=K'|X|=K'|X+C|$.

For the inductive step, assume $|C|>1$, let $\gamma\in C$ and $C'=C\setminus\{\gamma\}$. Then
\[X+B+C=(X+B+C')\cup\bigl((X+B+\gamma)\setminus(Z+B+\gamma)\bigr)\]
where
\[Z=\{x\in X\mid x+B+\gamma\subseteq X+B+C'\}.\]
$Z\subseteq X$ so by minimality $|Z+B|\geq K'|Z|$. We have
\begin{align*}
|X+B+C|&\leq |X+B+C'|+|(X+B+\gamma)\setminus(Z+B+\gamma)|\\
&=|X+B+C'|+|X+B|-|Z+B|\\
&\leq K'|X+C'|+K'|X|-K'|Z|\\
&=K'(|X+C'|+|X|-|Z|).
\end{align*}
Now we want to understand the right hand side $X+C$. Note that
\[X+C=(X+C')\sqcup\bigl((X+\gamma)\setminus(W+\gamma)\bigr)\]
where
\[W=\{x\in X\mid x+\gamma\in X+C'\}.\]
In particular this is a disjoint union, so
\[|X+C|=|X+C'|+|X|-|W|.\]
We also have $W\subseteq Z$ because $x+\gamma\in X+C'$ implies $x+B+\gamma\subseteq X+B+C'$. Thus $|W|\leq|Z|$, so
\[|X+C|\geq |X+C'|+|X|-|Z|,\]
which, when combined with the above inequality, completes the induction.
\end{proof}
\subsection*{整理说明（非作者正文）}
题目使用非空有限集以及非负整数 $m,n$，并取 $0B=\{0\}$；这是对原文使用范围和记号的补明。保留先用关键引理推出主定理、再证明引理的原顺序。开头记号段据同章定义整理，并非原文逐字引文。删去历史评论和图7.2的辅助解释；主证明不依赖该图。允许的标准先修只有有限集计数、阿贝尔群运算和归纳；最小扩张率子集的选择、$Z,W$ 构造、迭代和三角不等式均已收入。本题难点是用最小扩张率子集把单次增长控制传递到任意迭代，而不是给定 Pl\"unnecke 界后的直接应用。PDF 正文已取得，页面截图通道未返回图像；保存的是明确标示的 TeX 转录稿。
\subsection*{可修改的简短标注（非作者正文）}
$c$：小倍增集合的迭代增长。

$a$：最小扩张率子集；平移并集的差集分解；Ruzsa 三角不等式的单射；归纳迭代。

\texttt{cross\_theory\_status = yes}。由集合加法的基数增长选择最小扩张率子集，利用平移覆盖关系保持该增长率，再用可逆编码的单射转成迭代和差集界。位置：Lemma 7.25 的 $Z,W$ 构造与 Theorem 7.24 的回代。
标签为非穷尽、可修改初稿。
\subsection*{许可与修改记录}
原作者及作品见来源栏；来源采用 CC BY-NC-SA 4.0。本题证摘录和附加整理沿用该许可。2026-09-28：增加题号、来源定位、独立排版、整理说明和初稿标注；数学内容的取材范围及排版变化见整理说明。
\end{document}
